\chapter{Anexos}

	\begin{figure}[H]
		\centering
		\includegraphics[width=0.8\linewidth]{Images/Anexos/en-construccion1.png}
	\end{figure}


\section{Hard Processor System (HPS)}\label{hpssec}

    El HPS es un sistema que incluye un procesador quad-core ARM Cortex-A53, el cual permite a los usuarios migrar los dise\n os existentes desde dispositivos Intel Stratix 10 SoCs hacia dispositivos Intel Agilex SoC.

    Adem\'as, el HPS habilita las capacidades de virtualizaci\'on de hardware a nivel de sistema mediante la incorporaci\'on de una unidad de gesti\'on de memoria del sistema. Estas mejoras de arquitectura aseguran que los SoCs cumplan los requerimientos de los mercados actuales y futuros de sistemas embebidos, incluyendo comunicaciones al\'ambricas e inal\'ambricas, aceleraci\'on de centros de datos, y numerosas aplicaciones militares. %\cite{hps}

    \begin{figure}[H]
        \centering
        \includegraphics[width=0.8\linewidth]{Images/Anexos/hpsblck.png}
        \caption{Diagrama de bloques del HPS.}
        % \label{fig:enter-label}
    \end{figure}

\begin{longtable}{p{0.3\textwidth} | p{0.6\textwidth}}

\toprule
\multicolumn{1}{c}{\textbf{Caracter\'istica}} & \multicolumn{1}{c}{\textbf{Descripción}} \\
\toprule
\endfirsthead

\toprule
\multicolumn{1}{c}{\textbf{Caracter\'istica}} & \multicolumn{1}{c}{\textbf{Descripción}} \\
\toprule
\endhead

Unidad de procesamiento quad-core ARM Cortex-A53 MPCore &
\begin{itemize}
  \item 2.3 $MIPS/MHz$ eficiencia de instrucciones
  \item Frecuencia del CPU hasta 1.4 $GHz$
  \item Rendimiento total a 14. $GHz$ de 13,800 MIPS
  \item Arquitectura ARMv8-A
  \item Ejecuta instrucciones ARM de 64 y 32 bits
  \item Instrucciones Thumb 16 y 32 bits para una reducci\'on del 30\% de uso de memoria
  \item Arquitectura de ejecuci\'on Jazelle RCT con bytecodes Java de 8 bits
  \item Pipeline fuera de orden, s\'uperescalar, longitud de variable, con predicción de saltos din\'amica
  \item Motor ARM NEON para procesamiento mejorado
  \item Unidad de punto flotante de precisi\'on simple y doble
  \item Tecnolog\'ia de de depurado y trazado CoreSight
\end{itemize} \\ \midrule

Unidad de gestión de memoria del sistema &

 Habilita un modelo de memoria unificada y extiende la virtualización de hardware a los periféricos de f\'abrica implementados en la FPGA \\ \midrule

Unidad de coherencia de caché &
 Cambios en los datos compartidos almacenados en cach\'e se propagan a trav\'es del sistema proveyendo coherencia bidireccional para elementos coprocesadores \\ \midrule

Memoria caché &
\begin{itemize}
  \item L1 Cache\begin{itemize}
        \item Cach\'e de instrucci\'on de 32 KB con verificaci\'on de paridad
        \item Cach\'e de datos de L1 de 32 KB con ECC
        \item Verificaci\'on de paridad
    \end{itemize}
  \item L2 Cache\begin{itemize}
      \item 1 MB compartido
      \item 8 vías asociativas
      \item Protección SEU con verificaci\'on de paridad en la memoria TAG y ECC en la memoria de datos RAM
    \item Soporte de blqueo de cach\'e
    \end{itemize}
\end{itemize} \\ \midrule

On-Chip Memory &
256 KB de memoria scratch interna \\ \midrule

Interfaces de memoria para el HPS externas SDRAM y memoria Flash &
\begin{itemize}
  \item Controlador de memoria dura con soporte para memoria DDR4 \begin{itemize}
      \item 40 bits (32 bits + 8 bits ECC) con selecci\'on de paquetes soportando hasta 72 bits (64 bits + 8 bits ECC)
      \item Soporte para memoria de hasta 32200 Mbps 
      \item Soporte ECC incluyendo calculos, correcci\'on de errores, correcci\'on de escritura diferida y contador de errores
      \item Planificaci\'on por prioridad configurable por sofware de rafagas de memoria SDRAM individuales
      \item Soporte de par\'ametros de tiempo completamente programables para todos los par\'ametros de tiempo especificos JEDEC
      \item Interfaz planificadora multipuerto (MPFE) hacia controlador de memoria integrado, que soporta AXI Quality of Service (QoS) para la interconexi\'on con la FPGA \end{itemize}
\item Controlador de almacenamiento para NAND flash\begin{itemize}
        \item asdfsdfasdf  
    \end{itemize}
\end{itemize} \\

Controladores de almacenamiento &
\begin{itemize}
  \item NAND Flash (ONFI 1.0, ECC por hardware)
  \item SD/MMC/eMMC 4.5, hasta 50 MHz, con DMA
  \item Soporte CE-ATA y comandos digitales
\end{itemize} \\

DMA &
\begin{itemize}
  \item 8 canales
  \item Hasta 32 interfaces handshake
\end{itemize} \\

Interfaces de comunicación &
\begin{itemize}
  \item 3x Ethernet 10/100/1000 (MAC con DMA)
  \item Soporte RGMII, RMII, GMII, MII, SGMII
  \item Compatible con IEEE 1588, VLAN, AVB
  \item 2x USB OTG (host/dispositivo, hasta 480 Mbps)
  \item 5x I2C, 2x UART 16550, 4x SPI (2M/2S)
\end{itemize} \\

Temporizadores e I/O &
\begin{itemize}
  \item 4 temporizadores generales + 4 watchdog
  \item 48 pines I/O directos
  \item Hasta 2 bancos IO96 asignables a DDR HPS
\end{itemize} \\

Puentes con FPGA &
\begin{itemize}
  \item HPS-to-FPGA AXI (32/64/128 bits)
  \item Lightweight bridge (AXI 32 bits)
  \item FPGA-to-SoC bridge (hasta 512 bits ACE-Lite)
  \item HPS-to-SDM / SDM-to-HPS bridge
\end{itemize} \\
\hline

\caption{Resumen de caracter\'isticas}
\end{longtable}

\section{Proyectos Personalizados para el Kit de Desarrollo}

\subsection{Configuraciones SmartVID en el Archivo QSF de Quartus Prime}

    El Agilex 7 ensamblado en este kit de desarrollo habilita la funci\'on SmartVID de forma predeterminada. Para evitar un error generado por Quartus Prime debido a configuraciones incompletas de SmartVID, se debe colocar las restricciones que se indican a continuaci\'on en el archivo QSF del proyecto de Quartus Prime. Dichas restricciones est\'an dise\~nadas para el circuito principal Intel Enpirion ED8401.

\subsubsection{DK-SI-AGI027FA y DK-SI-AGI027FC (No Intel Enpirion)}

    Abrir el archivo QSF del proyecto en Quartus Prime, copiar y pegar los siguientes scripts de restricci\'on en el archivo. Se debe asegurar de que no existan otras configuraciones similares con valores diferentes.

    \begin{lstlisting}
set_global_assignment -name USE_PWRMGT_SCL SDM_IO0
set_global_assignment -name USE_PWRMGT_SDA SDM_IO12
set_global_assignment -name USE_CONF_DONE SDM_IO16
set_global_assignment -name VID_OPERATION_MODE "PMBUS MASTER"
set_global_assignment -name PWRMGT_BUS_SPEED_MODE "100 KHZ"
set_global_assignment -name PWRMGT_SLAVE_DEVICE_TYPE LTC3888
set_global_assignment -name NUMBER_OF_SLAVE_DEVICE 1
set_global_assignment -name PWRMGT_SLAVE_DEVICE0_ADDRESS 62
set_global_assignment -name PWRMGT_VOLTAGE_OUTPUT_FORMAT "LINEAR FORMAT"
set_global_assignment -name PWRMGT_LINEAR_FORMAT_N "-12"
set_global_assignment -name PWRMGT_TRANSLATED_VOLTAGE_VALUE_UNIT VOLTS
    \end{lstlisting}

\subsubsection{DK-SI-AGI027FB y DK-SI-AGI027FES (Intel Enpirion)}

    Abrir el archivo QSF del proyecto en Quartus Prime, copiar y pegar los siguietes scripts de restricci\'on en el archivo. Se debe asegurar de que no existan otras configuraciones similares con valores diferentes.

    \begin{lstlisting}
set_global_assignment -name USE_PWRMGT_SCL SDM_IO0
set_global_assignment -name USE_PWRMGT_SDA SDM_IO12
set_global_assignment -name USE_CONF_DONE SDM_IO16
set_global_assignment -name VID_OPERATION_MODE "PMBUS MASTER"
set_global_assignment -name PWRMGT_BUS_SPEED_MODE "100 KHZ"
set_global_assignment -name PWRMGT_SLAVE_DEVICE_TYPE ED8401
set_global_assignment -name NUMBER_OF_SLAVE_DEVICE 1
set_global_assignment -name PWRMGT_SLAVE_DEVICE0_ADDRESS 62
set_global_assignment -name PWRMGT_VOLTAGE_OUTPUT_FORMAT "LINEAR FORMAT"
set_global_assignment -name PWRMGT_LINEAR_FORMAT_N "-13"
set_global_assignment -name PWRMGT_TRANSLATED_VOLTAGE_VALUE_UNIT VOLTS
    \end{lstlisting}

\subsection{Golden Top}

    Se puede utilizar el proyecto Golden Top como punto de partida para dise\~nos propios. Este proyecto ya viene cargado de restricciones, ubicaciones de pines, est\'andares de E/S definidos, direcciones y terminaciones generales. Las configuraciones de terminaci\'on de pines para DDR4 no est\'an incluidas.
